%!TEX root = /Users/markelikalderon/Documents/Git/perceptual_self-consciousness/perceptual_self-consciousness.tex
\chapter{Acknowledgments} % (fold)
\label{cha:acknowledgments}

During his brief, all too brief, tenure as a colleague at UCL, Jake McNulty and I read Sartre's \emph{La Transcendance de l'ego: Esquisse d'une description phénomenologique} and \emph{L'Imaginaire: Psychologie ph{\'e}nom{\'e}nologique de l'imagination} together. Our discussions in the reading group were invaluable to my thinking, and I have benefitted, as well, from subsequent discussions after his departure. While I had not yet begun to compose the present essay, I was keen to understand Sartre's notion of pre-reflective or non-positional consciousness. Our discussions were instrumental in transforming this initial curiosity into full-fledged commitment to the present undertaking. 

Since her arrival at UCL I have discussed ancient philosophy with Simona Aimar. A seminar that we co-taught on Aristotle's \emph{Metaphysica} {\sbl Θ} was invaluable for my thinking about perceptual powers generally, and I have learned a great deal from her paper, co-written with Carlotta Pavese, on \emph{technē} \citep{Aimar:2025tg}.

I am very grateful to have had the opportunity to offer seminars on Plato's \emph{Charmides} and Aristotle's \emph{De anima} 3.2, the topics of chapters~\ref{cha:offering} and \ref{cha:perceiving}, respectively. Teaching these seminars were instrumental to generating drafts of these chapters, and I am very grateful to all who participated.

% end acknowledgments (end)